% ============================================================
%  preamble.tex — packages, palette, listings, boites, macros
%  Charge depuis main.tex via \input
% ============================================================

% ---- encodage et langue ----
\usepackage[T1]{fontenc}
\usepackage[utf8]{inputenc}
\usepackage{lmodern}
\usepackage[french]{babel}
\usepackage{geometry}
\geometry{margin=2.5cm}

% ---- math, graphismes, liens, tables ----
\usepackage{amsmath, amssymb, amsthm}
\usepackage{graphicx}
\usepackage{hyperref}
\usepackage{xcolor}
\usepackage{listings}
\usepackage{booktabs}
\usepackage{array}
\usepackage{enumitem}
\usepackage{titlesec}
\usepackage{fancyhdr}
\usepackage{tcolorbox}
\usepackage{float}
\usepackage{microtype}
\usepackage{parskip}
\usepackage{tikz}

% ---- palette ----
\definecolor{codebg}{rgb}{0.96,0.96,0.96}
\definecolor{codeframe}{rgb}{0.75,0.75,0.75}
\definecolor{primary}{RGB}{0,90,160}
\definecolor{accent}{RGB}{200,60,0}
\definecolor{notegreen}{RGB}{0,120,60}

% ---- listings : python (par defaut) et benzai (texte brut, sans coloration) ----
\lstdefinestyle{python}{
  language=Python,
  backgroundcolor=\color{codebg},
  frame=single,
  rulecolor=\color{codeframe},
  basicstyle=\ttfamily\small,
  keywordstyle=\bfseries\color{primary},
  commentstyle=\itshape\color{notegreen},
  stringstyle=\color{accent},
  numbers=left,
  numberstyle=\tiny\color{gray},
  numbersep=8pt,
  breaklines=true,
  showstringspaces=false,
  tabsize=4,
}
\lstdefinestyle{benzai}{
  backgroundcolor=\color{codebg},
  frame=single,
  rulecolor=\color{codeframe},
  basicstyle=\ttfamily\small,
  breaklines=true,
  showstringspaces=false,
}
\lstset{style=python}

% ---- boites : Remarque (bleu) et Attention (orange) ----
\tcbuselibrary{skins, breakable}
\newtcolorbox{remarque}[1][]{
  colback=blue!5, colframe=primary, fonttitle=\bfseries,
  title=Remarque, #1
}
\newtcolorbox{attention}[1][]{
  colback=orange!8, colframe=accent, fonttitle=\bfseries,
  title=Attention, #1
}

% ---- theoremes : definition numerotee par section ----
\theoremstyle{definition}
\newtheorem{definition}{Définition}[section]

% ---- en-tetes / pieds de page ----
\pagestyle{fancy}
\fancyhf{}
\rhead{\textcolor{gray}{\small Implémentation CSP — Stage AMU M2}}
\lhead{\textcolor{gray}{\small \leftmark}}
\cfoot{\thepage}

\hypersetup{
  colorlinks=true,
  linkcolor=primary,
  urlcolor=primary,
  citecolor=primary
}

% ---- raccourcis math ----
\newcommand{\GD}{G_{\!D}}
\newcommand{\Tn}{\mathcal{T}(n)}

% ---- marqueurs verdicts (substituts aux emojis Unicode hors-BMP qui ne
%      passent pas avec inputenc utf8 selon le compilateur) ----
\definecolor{verdictplane}{RGB}{26,127,55}
\definecolor{verdictnonplane}{RGB}{34,34,34}
\newcommand{\bucketplane}{\tikz[baseline=-0.5ex]\filldraw[verdictplane] (0,0) circle (0.45ex);}
\newcommand{\bucketnonplane}{\tikz[baseline=-0.5ex]\filldraw[verdictnonplane] (0,0) circle (0.45ex);}

% ---- marqueur de date pour suivre l'evolution du document ----
\newcommand{\dateblock}[1]{%
  \par\medskip
  \begingroup
    \setlength{\fboxsep}{4pt}%
    \fcolorbox{primary}{blue!5}{%
      \color{primary}\footnotesize\bfseries
      \textsc{Ajout du #1}%
    }%
  \endgroup
  \par\smallskip
}
